\section{Solution Algorithm}
\label{sec:solution_algorithm}

We develop a hybrid solution approach combining exact methods for validation and metaheuristics for realistic-scale instances. This multi-method strategy ensures solution quality while maintaining computational tractability.

\subsection{Algorithm Design Philosophy}

Our algorithm design follows three principles:

\textbf{Principle 1: Progressive Solution Strategy.}
\begin{enumerate}
\item Solve deterministic Model 0 to obtain baseline routes
\item Use Model 0 solution as warm-start for Model 1 (stochastic)
\item Iteratively improve via local search
\end{enumerate}

\textbf{Principle 2: Scenario-Based Sampling.}
\begin{itemize}
\item Use $S = 100$ scenarios for reliability estimation
\item Validate with $S_{\text{val}} = 1000$ scenarios for accuracy
\item Statistical convergence criterion: relative error $< 2\%$
\end{itemize}

\textbf{Principle 3: Multi-Scale Hybrid Approach.}
\begin{itemize}
\item Small instances ($N \leq 20$): Exact methods (CPLEX)
\item Medium instances ($20 < N \leq 50$): Constructive heuristics + local search
\item Large instances ($N > 50$): Adaptive Large Neighborhood Search (ALNS)
\end{itemize}

\subsection{Algorithm 1: Adaptive Large Neighborhood Search (ALNS)}

ALNS \cite{Ropke2006} dynamically adapts search operators based on historical performance.

\subsubsection{Algorithm Structure}

\begin{algorithm}[H]
\caption{Adaptive Large Neighborhood Search for SVRP}
\begin{algorithmic}[1]
\STATE \textbf{Input:} Initial solution $x_0$, scenario set $S$, reliability threshold $\alpha$
\STATE $x \leftarrow x_0$, $x^* \leftarrow x_0$
\FOR{iteration $iter = 1$ to $maxIter$}
\STATE Select destroy operator $d$ with probability $p_d$
\STATE Select repair operator $r$ with probability $p_r$
\STATE $x' \leftarrow \text{Repair}(\text{Destroy}(x, d), r)$
\IF{$x'$ feasible and $f(x') < f(x)$}
\STATE $x \leftarrow x'$
\IF{$f(x') < f(x^*)$}
\STATE $x^* \leftarrow x'$
\ENDIF
\ENDIF
\STATE Update operator selection probabilities $p_d, p_r$
\ENDFOR
\STATE \textbf{Output:} Best solution $x^*$
\end{algorithmic}
\end{algorithm}

\subsubsection{Destroy Operators}

We implement 5 destroy operators:
\begin{enumerate}
\item \textit{Random Removal:} Remove $q$ random customers (default: $q = 0.1N$)
\item \textit{Worst Removal:} Remove customers contributing most to objective
\item \textit{Related Removal:} Remove geographically clustered customers
\item \textit{Route Removal:} Remove entire routes with lowest reliability
\item \textit{Time Window Violation Removal:} Remove customers causing violations
\end{enumerate}

\subsubsection{Repair Operators}

We implement 4 repair operators:
\begin{enumerate}
\item \textit{Greedy Insertion:} Insert customers at minimum cost position
\item \textit{Regret Insertion:} Insert customer with highest regret \cite{Potvin1996}
\item \textit{Reliability-Weighted Insertion:} Prioritize high-reliability insertions
\item \textit{2-Opt Exchange:} Perform 2-opt local search after insertion
\end{enumerate}

\subsubsection{Adaptive Mechanism}

For each operator $\omega$, we maintain:
\begin{itemize}
\item Usage count $u_\omega$
\item Score $s_\omega$ (updated based on solution quality)
\end{itemize}

Selection probability:
\begin{equation}
p_\omega = \frac{s_\omega}{\sum_{\omega'} s_{\omega'}}
\end{equation}

Score update:
\begin{equation}
s_\omega \leftarrow \begin{cases}
\rho_1 s_\omega & \text{if new global best} \\
\rho_2 s_\omega & \text{if better than current} \\
\rho_3 s_\omega & \text{if accepted} \\
\rho_4 s_\omega & \text{if rejected}
\end{cases}
\end{equation}

where $\rho_1 > \rho_2 > \rho_3 > \rho_4 = 1$. We use $\rho = (1.0, 0.9, 0.8, 0.7)$.

\subsection{Algorithm 2: L-Shaped Method for Two-Stage Model}

For Model 2 (two-stage stochastic program), we implement Benders decomposition (L-shaped method).

\subsubsection{Master Problem (First Stage)}

\begin{equation}
\min \quad c^T x + \theta
\end{equation}
subject to:
\begin{itemize}
\item Deterministic constraints
\item Optimality cuts: $\theta \geq \mathbb{E}[\pi^{(s)}] (h - Tx)$
\item Feasibility cuts (if infeasible)
\end{itemize}

\subsubsection{Subproblem (Second Stage) for Scenario $s$}

\begin{equation}
\mathcal{Q}(x, \xi^{(s)}) = \min \quad q^T y^{(s)}
\end{equation}
subject to:
\begin{equation}
Wy^{(s)} \geq h - Tx, \quad y^{(s)} \geq 0
\end{equation}

\subsubsection{Cut Generation}

Let $\pi^{(s)}$ be dual optimal of subproblem for scenario $s$. We add optimality cut:
\begin{equation}
\theta \geq \sum_{s \in S} p^{(s)} \pi^{(s)} (h - Tx)
\end{equation}

\textbf{Convergence:} Algorithm terminates when gap between upper bound (UB) and lower bound (LB) satisfies:
\begin{equation}
\frac{UB - LB}{UB} \leq \epsilon_{\text{gap}}
\end{equation}
with $\epsilon_{\text{gap}} = 0.01$ (1\% optimality gap).

\subsection{Algorithm 3: Clustering-Based Multi-Depot Algorithm}

For Model 3 (multi-depot SVRP), we use clustering-first, routing-second approach.

\subsubsection{Phase 1: Customer-Warehouse Assignment}

We use Reliability-Weighted K-Means:

\begin{algorithm}[H]
\caption{Reliability-Weighted Clustering}
\begin{algorithmic}[1]
\STATE \textbf{Input:} Customers $V$, warehouses $W$, reliability matrix $\alpha_{ij}$
\STATE Initialize: Assign each customer to nearest warehouse
\FOR{iteration $iter = 1$ to $maxIter$}
\STATE Calculate cluster centroids
\FOR{each customer $i \in V$}
\STATE Reassign to warehouse minimizing $\frac{d_{iw}}{\alpha_{iw}}$
\ENDFOR
\IF{no reassignments}
\STATE break
\ENDIF
\ENDFOR
\STATE \textbf{Output:} Clusters $\mathcal{C}_w$ for each warehouse $w \in W$
\end{algorithmic}
\end{algorithm}

\subsubsection{Phase 2: Independent Route Optimization}

For each warehouse $w$, solve independent SVRP on cluster $\mathcal{C}_w$ using ALNS (Algorithm 1).

\subsubsection{Phase 3: Inter-Warehouse Coordination}

If load imbalance exists:
\begin{enumerate}
\item Identify over-loaded warehouse $w_1$ and under-loaded $w_2$
\item Transfer customers near boundary between $\mathcal{C}_{w_1}$ and $\mathcal{C}_{w_2}$
\item Re-optimize routes for affected warehouses
\item Repeat until balanced or max iterations reached
\end{enumerate}

\subsection{Algorithm 4: Multi-Objective Optimization}

For Model 4, we use the $\epsilon$-constraint method.

\subsubsection{Procedure}

\begin{enumerate}
\item Calculate Pareto optimal bounds for each objective
\item Discretize objective space into grid
\item For each grid point $(\epsilon_2, \epsilon_3)$:
\begin{itemize}
\item Solve: $\min f_1$ s.t. $f_2 \leq \epsilon_2, f_3 \leq \epsilon_3$
\item Record solution if feasible
\end{itemize}
\item Filter dominated solutions to obtain Pareto frontier
\end{enumerate}

\subsubsection{Visualization}

We visualize Pareto frontier in 3D objective space (time-reliability-cost). Decision-makers can select preferred trade-off point based on operational priorities.

\subsection{Implementation Details}

\subsubsection{Computational Environment}

\begin{itemize}
\item Language: Python 3.10
\item Optimization solver: CPLEX 12.10 (for exact methods)
\item Parallel computing: 8-core Intel i7 processor, scenario evaluations parallelized
\item Maximum runtime: 1 hour per instance (industrial practice standard)
\end{itemize}

\subsubsection{Convergence Criteria}

\begin{itemize}
\item ALNS: 10,000 iterations or no improvement for 1,000 iterations
\item L-Shaped: 1\% optimality gap or 100 Benders iterations
\item Multi-objective: 50 Pareto points or computational limit
\end{itemize}

\subsubsection{Validation}

We validate algorithms by:
\begin{enumerate}
\item Comparing heuristic solutions to CPLEX optimal solutions (small instances)
\item Testing multiple random seeds, reporting mean and standard deviation
\item Cross-validating on hold-out scenario set ($S_{\text{val}} = 1000$ scenarios)
\end{enumerate}

\subsection{Algorithm Performance Summary}

Table 2 summarizes algorithm performance (details in Section \ref{sec:results}):

\begin{table}[H]
\centering
\caption{Algorithm Performance on Test Instances}
\begin{tabular}{lcccc}
\toprule
Instance Size & Algorithm & Mean Time (sec) & Gap to Optimal & Std Dev \\
\midrule
Small ($N=15$) & CPLEX & 12.3 & 0.0\% & - \\
Medium ($N=30$) & ALNS & 45.7 & 2.3\% & 1.1\% \\
Large ($N=60$) & ALNS & 183.2 & 3.8\% & 1.9\% \\
Multi-Depot & Clustering+ALNS & 215.6 & 4.1\% & 2.3\% \\
\bottomrule
\end{tabular}
\end{table}

These results demonstrate that our hybrid approach achieves near-optimal solutions (within 5\% gap) for realistic instances within acceptable computational time.
